% 逐条回复信 —— IEEE Access 重投件 Access-2026-49387
% 编译：见同目录 build.sh（须在纯 ASCII 目录编译）
\documentclass[10pt,letterpaper]{article}
\usepackage[T1]{fontenc}
\usepackage[utf8]{inputenc}
\usepackage{mathptmx}
\usepackage[margin=1in]{geometry}
\usepackage{enumitem}
\usepackage{titlesec}
\usepackage{fancyhdr}
\usepackage{booktabs}
\usepackage{array}
\usepackage{longtable}
\usepackage{xcolor}
\usepackage[hidelinks]{hyperref}

\setlist{nosep,leftmargin=1.5em}
\titleformat{\section}{\normalfont\large\bfseries}{\thesection}{0.6em}{}
\titleformat{\subsection}{\normalfont\normalsize\bfseries}{\thesubsection}{0.6em}{}
\titlespacing*{\subsection}{0pt}{1.2ex plus .2ex}{0.6ex}
\pagestyle{fancy}
\fancyhf{}
\fancyhead[L]{\footnotesize Response to Reviewers --- Access-2026-49387}
\fancyhead[R]{\footnotesize\thepage}
\renewcommand{\headrulewidth}{0.4pt}

\newcommand{\Concern}[1]{\par\smallskip\noindent\textbf{Concern.}\ #1\par}
\newcommand{\Response}[1]{\par\smallskip\noindent\textbf{Response.}\ #1\par}
\newcommand{\Action}[1]{\par\smallskip\noindent\textbf{Action.}\ #1\par}
\newcommand{\Remains}[1]{\par\smallskip\noindent\textbf{What remains.}\ #1\par}

\begin{document}

\begin{center}
{\large\bfseries RESPONSE TO REVIEWERS}\\[0.6ex]
{\normalsize Manuscript \#Access-2026-49387}\\[0.4ex]
{\normalsize \textbf{Revised title:} \emph{Two Reachability-Gated Relations in MST-wi:
Temporal Mismatch, Oracle Admissibility, and Measurement Limits}}\\[0.4ex]
{\normalsize Qiulong Lin and Nanlei Chen}\\[0.4ex]
{\normalsize 19 October 2026}
\end{center}

\vspace{0.5ex}
\noindent Dear Associate Editor and Reviewers,

Thank you for the detailed reports. The revision is substantive rather than
cosmetic, and we have kept the narrowing that both reports asked for: the title,
the abstract, the running head and every claim are now restricted to the two
reachability-gated relations that were actually examined. Three statements that
were previously argued in prose are now computed from records and reproduced by
script---the complete factorial design and its reconciliation, the truth table of
the two relations, and a competitive differential baseline on a matched request
budget. One experiment that had not been run at all (a competing differential
baseline) has now been run and reported. Where a request could not be met with
the evidence we have, we say so in the manuscript instead of softening the text;
eight of the twenty-one items are met only in part, and Section~\ref{sec:partial}
of this letter lists them together with what is missing and why.

References below are to the revised manuscript (26 pages, 22 tables). Table and
section numbers are stable; page numbers are given for the tables most often at
issue. Every number quoted in this letter is generated by a script from the raw
records; none is typed by hand, and the generation scripts are part of the
reproducibility package.

\section{Reviewer 1}

\subsection{Novelty positioning against existing BOLA testing research}
\Concern{The paper identifies limitations in reachability-based adjudication,
oracle assumptions and metric design, but does not sufficiently distinguish these
contributions from established work on oracle problems, stateful web testing,
authorization testing and metamorphic testing. A structured comparison is
requested.}
\Response{We agree that ``a frozen snapshot goes stale'' is not new, and the
manuscript now says so explicitly. What we can support with evidence is narrower
and is stated as three falsifiable increments rather than as a general
positioning claim. (i)~\emph{Predicate level:} at source, the verdict of each of
the two relations is a function of the crawl-time reachability model and the
request outcome, so it is determined by quantities that a direct
request-plus-oracle comparison already has (Section~VII-B, confirmed by
execution in Section~VII-C). (ii)~\emph{Subset level:} the false proofs we
exhibit are confined to the cells whose grant is invisible to the framework's own
record, and the channel decomposition (Table~11, p.~13) shows exactly which
channels those are; neither product's membership surface produces a single false
\emph{allow}. (iii)~\emph{Denominator level:} a pooled input-level specificity
cannot resolve a failure confined to one class, and no fixed threshold can force
such a class below it (Proposition~1 and Corollary~1, Section~IV-E). We also
state plainly that we do not compare detection rates and cannot, from our
evidence, claim any.}
\Action{Section~IX rewritten around these three increments; the claim of being
``orthogonal'' to prior work deleted; the abstract's first sentence and the
conclusion restated in the same terms.}

\subsection{A complete definition of the false-proof outcomes}
\Concern{The mapping between alarms, false alarms and proven violations and the
ground-truth authorization state is not formalized consistently; decision
variables and truth conditions are requested for every combination of request
outcome, actual entitlement, oracle response and detector verdict, together with
an account of how a legitimate successful request is distinguished from
unauthorized access without letting the tested reachability model determine its
own ground truth.}
\Response{Table~16 (Section~VI-B, p.~17) is now that table. It defines four
quantities over the 712 temporal cells: $Y$ (the ordinary request succeeds), $M$
(the crawl-time model lists the object for the subject), $O$ (the introspection
surface's answer, with an explicit ``no definite answer'' value), $A$ (the alarm
of either relation; on this domain both reduce to $y\wedge\neg m$) and $D$ (the
direct adjudicator, $y\wedge\neg o$), plus $H=A\wedge\neg o$, the promoted verdict
of the state-calibrated method, which exists only where $O$ is definite. The
twelve rows cover every combination of $(Y,M,O)$, and the count columns split
each row by the ledger's entitlement $E$, so each row states its truth conditions
for both values of $E$. On the anti-circularity point: the label $E$ comes from
the framework's own grant and membership ledger, which is an independent
entitlement source; $M$ is an \emph{input} to the relations and is never used as
a label. A legitimate success is separated from unauthorized access only by an
independent \emph{deny} from $E$ together with a content-level \emph{allow}
observation---never by $M$---and the manuscript says that the relation's alarm
is not evidence of illegal access.}
\Action{Table~16 and the definitions in Section~IV-A; the three-way split of the
native run into $A$ (the engine's own alarm), $H$ (our promotion policy) and $T$
(confirmed) is given in Table~21 (p.~21).}

\subsection{Characterization of authorization-state transitions}
\Concern{The three-phase protocol does not comprehensively model revocation,
inherited permissions, public visibility, role changes and concurrent updates;
the state-transition semantics, the timing of each observation and the
consistency assumptions behind the deterministic false proofs are requested, as
is whether the conclusions hold when authorization changes between request and
adjudication.}
\Response{Section~V-A now states the semantics that the experiment actually
exercises, and three scope limits are stated as part of the result rather than
buried in the discussion. The eight committed transitions are a designed set
(direct grant and revoke, team attach and detach, member addition and removal,
private-to-public visibility, site-administrator grant and revoke); each cell is
serialised, no mutation is issued between a cell's request and its probe, so the
window between them is static \emph{by construction}; and a deployment in which a
grant or revoke lands inside that window is outside what these cells measure---for
such a case the cell's conclusion does not apply and the observation must be
repeated, which the manuscript says we do not claim to detect. Revocation and
inheritance behave differently from grants in this design and the table shows it:
the revoke rows produce no stale false proof at all (Table~12, p.~14), because an
over-listing model makes the relation silent (a miss) rather than wrong (a false
proof). On concurrency, the harness is single-threaded, no concurrent writer ran,
and $hv$ is declared to be a harness-maintained logical version rather than an
atomic server snapshot---which is what makes the still window meaningful.}
\Action{Section~V-A (Table~12, p.~14); the same scope limits restated in
Section~VIII.}

\subsection{Operationalizing the admissibility condition}
\Concern{Clause~(2)---that the oracle is not affected by the same defect under
investigation---remains a qualitative judgement; a reproducible procedure for
determining admissibility \emph{before} examining outcomes is requested.}
\Response{We agree this is the weakest link, and we have not solved it. What we
did do: the candidate surfaces, the granting channels and the defect classes were
fixed in a protocol file whose digest is bound into the run record before the
reported run; the primary-vantage gate was checked (4 of 4) before that run; and
the one amendment made afterwards is bound by digest to the run record, with the
contradiction in its own timestamp disclosed rather than back-dated---the
amendment file's \texttt{written\_at} postdates the last probe, and the
manuscript says so and gives the commands to re-check it. What we did \emph{not}
do is supply a mechanical, implementation-independent test of independence from
the tested mechanism. Section~VI-E now says explicitly that clause~(2) is a
per-system judgement and is not automatable from readability alone, and it makes
the judgement concrete instead of abstract: Gogs and GitLab fail it, and the
failure is observable as false proofs on identifiable channels (Tables~11
and~13). We would rather state the limit than imply a procedure we do not have.}
\Action{Section~VI-E; the pre-registered protocol and the amendment timing note
are in the reproducibility package.}
\Remains{No mechanical independence criterion. This is listed as a limitation in
Section~VIII.}

\subsection{Proposition 2: separating the definitional from the empirical}
\Concern{The proposition risks being tautological; the formal identity should be
distinguished from the empirical premise that the evaluated relations consume no
additional decision-relevant information, with a truth-table derivation or
explicit functional representation.}
\Response{The result is now split in two. \textbf{Lemma~1} (absorption) is the
algebraic part: if a verdict is a function of an oracle answer and an ordinary
request, it cannot outperform the direct consumer of those same inputs. That part
is definitional and is labelled as such. \textbf{Proposition~2} is the empirical
part: the two relations fall inside the lemma's domain, which is a claim about
this implementation and is checked, not assumed. The explicit predicates arewritten out ($A_{002}=\neg sup(B,A)\wedge\neg m\wedge\neg err(o_1)\wedge
q(o_1,o_2)$ and $A_{004}=y\wedge\neg m$, against $D=y\wedge\neg o$), and on the
domain studied both reduce to $y\wedge\neg m$. The reduction was verified
mechanically against 576 constructed records and 712 temporal cells with zero
mismatches; the generating script exits rather than emitting the table if any
record disagrees. The equivalence is stated correctly as $A=D \Leftrightarrow
y=0 \vee m=o$: where $D$ is defined the two verdicts coincide in 530 of the 584
readable cells and part in the two rows with $Y{=}$allow and $M\neq O$---40 cells
where the model is right and the oracle wrong, and 14 where the model is stale
and the oracle right. Those two rows are the two failure mechanisms the
cross-vendor section keeps apart. The falsification condition is stated: produce
an observable, outside what a succeeding request plus a readable oracle can
determine, on which an authorization metamorphic relation still decides.}
\Action{Section~VI-B with Lemma~1, Proposition~2, Table~16 (p.~17) and the
corrected equivalence; the earlier claim of extensional equality is deleted.}

\subsection{Coverage and identifiability of the 576-record enumeration}
\Concern{A transparent mapping from configurations and scenarios to all records
is requested, with the complete factorial design, an account of how each factor is
varied, attribution of the reported effects, and an explicit reconciliation of the
36/18 representative grid with the later 312/264 partition.}
\Response{The enumeration is now inspectable end to end. Table~1 (p.~4) is the
complete design: seven factors, their levels, how each level is varied, the record
count per level and what it crosses with, with the additive identity
$44\cdot2\cdot3\cdot2+8\cdot3\cdot2=576$ in the footnote. Table~19 (p.~20) is the
configuration-by-scenario census, where a dash marks an inadmissible combination
and the caption says why (a configuration with no share has nothing to list; a
listed cell requires introspection on). The reconciliation the reviewer asks for
is Table~5 (pp.~8--9): it starts from the S3--S7 original grid of 54 records
(which contains 8 v2 false proofs) and adds the remaining blocks, asserting by
record key that the blocks are mutually disjoint and that their union is exactly
576, split 312 readable / 264 unreadable. Section~III-A states what the design
can and cannot identify: the three dimensions are nested rather than crossed (the
relation is single-valued on S5; channel and visibility are defined only inside
one exception setting), so a full interaction estimate is not available---which is
why attribution is done with matched pairs (Table~2, p.~5) rather than with a
regression.}
\Action{Tables~1, 2, 5 and 19; Section~III-A.}

\subsection{Comparison against competitive black-box methods}
\Concern{\texttt{direct\_only} is deliberately unfair and \texttt{obs\_success} and
\texttt{always\_alarm} are intentionally weak; a differential, multi-account
authorization tester probing equivalent objects under distinct privilege levels
is requested, with matched query budgets, execution costs and detection
outcomes.}
\Response{This is the one comment that required a new experiment, and it has been
run. Section~VI-D and Table~17 (p.~19) report three conventional differential
rules on the same 42 controlled cells (six configurations crossed with S1--S7)
and, critically, on the \emph{same probe pair}: every method consumes one request
as the owner and one as the actor; ours issues one additional read of the
authorization surface on the 27 cells where a relation alarms and none where it
does not (0.64 per cell on average). \emph{B1 status-only} alarms whenever the
replayed request returns a non-error. \emph{B2 equality replay} alarms only when
the owner succeeded and the actor received byte-identical output---the
differential rule with no authorization model. \emph{B3} adds the strongest
information a tester can hold without reading the target: it suppresses every
cell covered by the tester's own share log. On these cells B1 and B2 claim on all
27 allowed cells with 12 false proofs each (precision 55.56\%), B3 claims 21 with
6 false proofs (71.43\%), all six on the out-of-band channel, and ours claims 10
with no false proof but 17 abstentions (coverage 59.52\%). Two readings are
explicitly ruled out in the text: this is \emph{not} a detection gain, because
the baselines find all fifteen violations while ours confirms the ten whose
surface is readable and declines the other five; and it is \emph{not} a free
improvement, because coverage falls to 59.52\%. A corroborating observation on a
real system is reported with the same qualification: on the Gitea snapshot the
equality rule fires on all 25 legitimately allowed cells and the in-band log
still leaves 24, because the framework log does not observe 0.96 of the effective
access there---a property of the constructed topology, not an estimate of field
prevalence.}
\Action{Section~VI-D and Table~17 (p.~19); the runner
(\texttt{experiments/run\_baseline.py}) and its raw output are in the
reproducibility package, and both baseline tables are reproduced byte-identically
by \texttt{verify\_artifact.sh repro}.}
\Remains{Execution cost is reported as request counts, not wall-clock time, for
these four methods: all of them replay the same recorded tape, so a timing
comparison would measure our replay loop. End-to-end timing is reported only for
the real-system temporal loops (Gitea 21.2~s, GitLab 7.3~s, of which the probes
are 20.7~s and 6.0~s).}

\subsection{Channel-by-channel isolation of the oracle-validity problem}
\Concern{A channel-by-channel confusion matrix separating direct membership,
team-derived access, administrative grants, public visibility and out-of-band
sharing is requested, together with a distinction between an incorrect oracle
response and a stale reachability model.}
\Response{Table~11 (Section~V-B, p.~13) is that matrix: the owner-vantage
introspection surface scored against the ledger per channel and per product, with
cells, agreements in each direction, the two error directions, and a separate
``unanswered'' column for anonymous readers. Neither product produces a single
false \emph{allow}: the membership surface errs only by omission, and public and
site-administrator cells are exactly the channels that are not group or
repository membership. The two failure mechanisms are kept apart by construction
and by table: in the truth table, row~6 is the oracle answering \emph{deny} on a
cell the ledger allows (40 cells, an unsound oracle), row~7 is the model failing
to list a cell the ledger allows (14 cells, a stale model), and the two verdicts
disagree in precisely those rows. On Gogs: the identical three-phase protocol was
run (Table~13, p.~14), but Gogs' introspection surface exposes no per-channel
decomposition equivalent to the other two products, so the manuscript declares
the gap instead of manufacturing the matrix.}
\Action{Tables~11 and~13; the gap stated in Section~VIII.}
\Remains{No per-channel four-cell matrix for Gogs.}

\subsection{Detection results across independently implemented defects}
\Concern{The three seeded instances share one fault pattern; the positive control
should be expanded to multiple independently implemented authorization defects,
including unauthorized reads, unauthorized modifications, inherited-permission
errors and privilege escalation, with per-defect recall and false-positive
rates.}
\Response{We did not implement multiple independent defects, and the manuscript
says so. What was added is a real vulnerable/fixed pair rather than another
seeded instance: a publicly fixed defect in Gitea (object-versus-path ownership
in the comment-permission check, fixed between 1.20.5 and 1.20.6), driven through
the same adjudication gate. The verdict uses the ledger's entitlement and the
content-level observation: on 1.20.5 the gate yields one true positive and no
false proof; on 1.20.6 it yields no alarm and no claim; the twelve control checks
(ten expected responses across the two builds and two recovery checks) all behave
as expected, and the generating script refuses to emit a number if any control
deviates. The manuscript is explicit that this is a positive control of the
\emph{gate}---an independent entitlement deny together with a content-level
allow---and not evidence about the relations' recall on real defects, and the
conclusion states that a stronger claim needs a \emph{population} of real
defects, which a single pair cannot provide. The seeded control is also bounded
honestly: the three instances instantiate one fault pattern, and the one BFLA
defect we seeded (a non-privileged actor succeeding at a collaborator-write
endpoint) is reported as a boundary of the read-oriented logic, not as a
detection failure.}
\Action{Section~V-D and Table~15 (p.~15); Section~V-C and Table~14 (p.~14) for
the seeded control; the scope statement in Section~X.}
\Remains{Per-defect recall and false-positive rates across several independent
defect classes. Not done, and stated as out of scope in the manuscript.}

\subsection{Control of temporal and observational effects}
\Concern{The ordering of crawling, entitlement mutation, oracle querying and
request execution must be specified, with timestamps or event logs sufficient to
verify that the comparison changes only the model-reading time, and an account of
how live-state queries avoid changing the tested state or gaining privileges.}
\Response{The ordering is fixed and stated: crawl, then freeze the model at
$t_0$, then commit the transition, then query the oracle, then issue the request;
all observations are \texttt{GET} requests, so nothing in the measurement path
mutates state; the recrawl uses the same credential as the original crawl, so it
grants no privilege the detector did not have; and the owner oracle is reported
separately from the subject-side reading rather than pooled with it. The run
record carries the protocol digest, the event sequence numbers, the policy
version, the $M_0$/$M_1$ model digests and the per-request timestamps, and the
harness version $hv$ is declared to be a logical version rather than an atomic
server snapshot. The results are symmetric across the two products: Gitea 280
cells with 13 stale false proofs, all 13 repaired by a recrawl with the same
credential and 0 mismatches under the live reading; GitLab CE 432 cells with 16
stale false proofs, all 16 repaired and 0 under live reading.}
\Action{Section~V-A and Table~12 (p.~14); the pre-registered protocol file, the
run records and the amendment timing note are in the reproducibility package.}

\subsection{The specificity denominator}
\Concern{The 84.62\% figure assumes two follow-up inputs per case and cannot
establish the numerical behaviour of the published 99.81\% figure; the original
counts or a parameterized sensitivity analysis are requested, with a clear
distinction from conventional specificity.}
\Response{The published figure is now treated as what it is: a specificity,
$\mathrm{TN}/(\mathrm{TN}+\mathrm{FP})$, over the follow-up inputs $F$, with the
manual verification of the reported failures acknowledged---so the labels are not
the crawler's own answer, and we do not attribute the number to a circular label.
The 84.62\% illustration has been removed from the manuscript entirely, together
with the claim that it agrees with the crawler. We cannot recover the engine's
exact input count, so instead of substituting a value we parameterize: with
$\mathrm{FN}=0$ the figure equals $1-\mathrm{FP}/(F-\mathrm{TP})$, a class
supplying a share $w$ of the negative inputs can move it only within $[1-w, 1]$,
and no fixed threshold $t<1$ can guarantee that a class-uniform failure pushes it
below $t$ (Proposition~1 and Corollary~1, Section~IV-E). The proposed
scenario-level quantity is labelled an analytical construct for characterising
resolution, not a replacement metric, and the text says the finding is about what
the published figure can be asked to show, not about its arithmetic.}
\Action{Section~IV-E; restated in Section~VI-C; the figure removed from the body.}

\subsection{Heterogeneous scenario weights}
\Concern{Proposition~1 assumes equal contributions per scenario; the result
should be extended to unequal request counts, correlated inputs and different
alarm counts per scenario, including the case where the number of generated
inputs depends on scenario complexity or detector behaviour.}
\Response{The fixed-weight case is derived (Proposition~1 and its corollary,
Section~IV-E) and the arbitrary-share case is bounded without any equal-weight
assumption: whatever the class shares are, a failing class contributing a share
$w$ of the negative inputs moves the figure only within $[1-w,1]$
(Section~VI-C). We did \emph{not} derive the endogenous case in which the number
of generated inputs is itself a function of detector behaviour, because our
records are a fixed enumeration in which the input count does not respond to the
detector; deriving it would require a pipeline whose generation is
detector-dependent, which we did not build. Rather than assert a generality we
cannot support, Section~VIII states the limitation: records within a design cell
are not independent, the records are a fixed enumeration, and we therefore report
counts and rates rather than inferential intervals, with no result resting on
statistical significance.}
\Action{Section~IV-E and Section~VI-C; the limitation stated in Section~VIII.}
\Remains{No derivation for detector-endogenous input counts.}

\subsection{Consistent treatment of abstentions and operational cost}
\Concern{Perfect precision bought by abstention is not characterized by precision
and recall alone; coverage, selective risk, abstention rate and performance
conditional on oracle availability are requested, alongside execution time and
request overhead, with a distinction between a safe abstention and a successful
detection.}
\Response{Table~4 (Section~IV-D, p.~8) rescored the same 576 records on the
\emph{decision} rather than on the alarm, and reports coverage, the abstention
count and rate split by whether the record held a true violation, false proofs,
misses and selective risk, with the explicit warning that risk is comparable
across methods only at equal coverage---which these methods do not have. The
state-calibrated method covers 60.42\% and abstains 228 times; panel~(b) splits
those abstentions by reason, and only the first row (156 surface-unreadable, 108
of them true violations) is the availability gate the method is named for, the
other 72 being a false proof converted into a non-decision. The direct methods
abstain under the same unreadable-means-abstain contract (coverage 54.17\%, 264
abstentions), and the forced-negative reference---coverage 100\%, 108 misses,
risk 18.75\%---is labelled as the price of forcing a verdict on unreadable cells,
not as a property of any method. A passed record is defined as one on which the
method did not fire, explicitly not a certificate of safety, and the text refuses
to rank by risk alone, noting that the always-degrade reference has the
second-lowest risk while detecting nothing. Cost is reported as request counts:
4833 enumeration requests (median 8, range 1--17) and 104 probe \texttt{GET}s
(52 cells $\times$ 2).}
\Action{Table~4 (p.~8) and Section~IV-D; the abstention semantics unified across
methods in Section~IV-A.}
\Remains{Wall-clock execution time for the replay-based methods, for the reason
given under comment~7.}

\subsection{Reproducibility of the principal claims}
\Concern{A version-pinned environment, exact commands, complete configuration and
seed files, machine-readable raw outputs and scripts that regenerate every table
and figure are requested, together with documentation of the two failed relation
executions, the engine-side defect and the treatment of missing or failed
observations.}
\Response{The package now carries a pinned environment, the exact commands, the
seed and configuration files, the machine-readable raw records and the scripts
that regenerate every table and figure, and it was verified from a clean copy
rather than asserted: in an empty directory the submission package compiles
(26 pages, no undefined references), and \texttt{verify\_artifact.sh repro}
reproduces all 28 generated artefacts byte-identically, including the two
baseline tables. The native run is reported as three stages rather than one
number (Section~VII-C, Tables~20--22, pp.~21--22): $A$, the engine's own
\texttt{getFailures()} alarms, 6; $H$, our promotion policy applied to them, 6;
$T$, confirmed, 0---with 6 false proofs. The content predicate is reported by
invocation count rather than claimed clean, and the text says so. The two failed
relation executions and the engine-side defect are documented at source, and the
treatment of missing or failed observations is written into the protocol file
that the run record binds.}
\Action{Section~VII-C and Tables~20--22; \texttt{reproduce.sh},
\texttt{verify\_artifact.sh} and the hash manifests in the package.}
\Remains{The failure and missing-observation rules live in the protocol file and
the scripts; Section~VII states them only in part.}

\subsection{Narrowing the conclusions; an evidence-to-claim table}
\Concern{Statements suggesting that adjudication of genuine violations is absent
from the evaluated relations risk extending beyond the experiments; conclusions
established by source-level dependence, by the synthetic construction and on real
systems should be distinguished, with a final table linking each research
question to its assumptions, tested configurations, results and limitations.}
\Response{Table~18 (Section~VI-E, p.~20) is that table. Each row names the
assumption the claim needs, the artefact that carries it, and the boundary that
remains; the last row is the adjudication gate and is marked as not being
evidence about detection effectiveness. The three findings are explicitly
attributed to different kinds of evidence---source-level functional dependence, a
constructed enumeration, and two real-system studies---and the text says a reader
should be able to see which kind carries which sentence. The sentence that
extended beyond the experiments has been deleted, together with the
mis-attribution of a constructed result to a published one; the title, abstract,
running head and conclusion are narrowed to the two reachability-gated relations,
and each finding is stated with its own falsification condition.}
\Action{Table~18 (p.~20); Section~VI-E; title, abstract and Section~X.}

\section{Reviewer 2}

\subsection{Novelty, and whether re-crawling or live reading removes the problem}
\Concern{State what is new beyond showing that a snapshot goes stale, and whether
re-crawling or live reading removes the problem; the data show 0 mismatches under
live reading on Gitea.}
\Response{Live reading removes the measured false proofs on both products: 0
mismatches under the live reading against 13 stale false proofs on Gitea and 16
on GitLab, and a recrawl with the same credential repairs 29 of 29. We removed
the unconditional claim that more crawling is ineffective, and the text now says
when a recrawl suffices---when the model can be refreshed with the same
credential before adjudication---and when it does not: a grant the surface cannot
enumerate at all is invisible to any recrawl, which is the out-of-band channel
where the residual false proofs sit. What is new is therefore not the observation
that a snapshot goes stale, but the measurement of when staleness produces a
\emph{false proof} rather than a miss, the separation of that mechanism from an
unsound oracle, and the fact that the residual survives the strongest
tester-side bookkeeping.}
\Action{Section~V-A (Table~12, p.~14); Section~VI-A; the related-work claim
restated in Section~IX.}

\subsection{Scope of claims: two relations of one family}
\Concern{Only 2 of 76 relations were studied and only 2 of 11 authorization
relations fired on Gitea; the title and abstract should be narrowed or the
evaluation extended.}
\Response{The title, the running head, the PDF title and the abstract are
narrowed to the two reachability-gated relations in MST-wi, and the abstract
states the family restriction and the native trigger range: of the 11 catalog
authorization relations, 10 were executed against Gitea and only the
reachability-gated ones alarm. We did not extend the evaluation to other relation
families; Section~VIII
names this as the paper's single most important limitation and says what raising
it would require---relations gated otherwise, which is a different experiment.}
\Action{Title, running head, abstract, Section~VIII.}

\subsection{Proposition 2}
\Concern{It is close to true by definition; separate what is empirical from what
is definitional.}
\Response{Done as described under comment~5: the algebraic absorption is a lemma
and is labelled definitional, the claim that these two relations fall in its
domain is a separate proposition with an empirical premise that is checked
against the records, and the truth table makes the functional representation
explicit.}
\Action{Section~VI-B, Table~16 (p.~17).}

\subsection{Detection side: no real system was vulnerable}
\Concern{Either add a real vulnerable case or remove the claim that the
metamorphic step adds nothing for detection.}
\Response{Both halves were done. A real vulnerable/fixed pair was added
(Section~V-D, Table~15, p.~15): a publicly fixed Gitea defect, one true positive
on the vulnerable build, no claim on the fixed build, twelve control checks
passing. And the claim was narrowed: the manuscript no longer says the
metamorphic step adds nothing \emph{for detection}; what it supports is that the
step adds no adjudication information where a readable, sound oracle exists
(Lemma~1 and Proposition~2), and that the relations' recall on real defects is
unevaluated---which Section~X states as the boundary of the evidence.}
\Action{Section~V-D and Table~15; Section~VI-B; Section~X.}

\subsection{The 84.62\% figure}
\Concern{It assumes two inputs per case and is only illustrative; keep it out of
the abstract and do not present it as quantitative support.}
\Response{The figure is absent from the abstract and has now been removed from
the body as well; it appears nowhere in the manuscript and is not used as
quantitative support for anything. The denominator argument is made in
parameterized form instead (Section~IV-E).}
\Action{Section~IV-E; verified absent from the compiled manuscript.}

\subsection{Constructed ground truth and the reported rates}
\Concern{Discuss how far the false-proof rates reflect the scenario design, such
as the per-cell multipliers, rather than the technique.}
\Response{The manuscript states in several places that the channel proportions
are topology design parameters and not field prevalence, that the eight
transitions are a designed set and not a sample of deployment changes, and that
the out-of-band share quoted on the real system is a property of the constructed
topology. The design dependence is also made inspectable rather than asserted:
Table~5 gives the per-block record counts behind the 312/264 partition, Table~19
gives the per-cell census, and the matched-pair attribution (Table~2) shows
exactly which cells carry each effect---the false-proof saving falls on one
scenario, the legitimate share created through the framework's own facility,
while the out-of-band scenario agrees with the reference. The reported rates are
therefore presented as properties of the construction, not as deployment
rates.}
\Action{Sections~V-B, VI-A and VIII; Tables~2, 5 and~19.}
\Remains{We did not recompute the constructed rates under equal per-cell
weights in parallel with the original weights; the sensitivity to the weighting
is described qualitatively rather than tabulated.}

\section{Answers to the additional questions}
\label{sec:addq}

\noindent\textbf{Reviewer 1.} (1)~Contribution to the body of knowledge: yes, in
the narrowed form stated in Section~IX and Table~18---a predicate-level account
of what the two relations determine, a channel decomposition of the false proofs,
and a resolution argument about a pooled specificity. (2)~Technically sound: the
algebraic part is separated from the empirical part and the empirical part is
checked against records by script; the unsupported parts are marked as
limitations. (3)~Comprehensive: the design, the ablation, the real-system studies
and the reproducibility package are all reported, with the gaps named. (4)~and
(5)~References: applicable and sufficient; none is inappropriate.

\noindent\textbf{Reviewer 2.} (1)~Contribution: partly, and we have narrowed the
claim accordingly rather than defending the broader reading. (2)~Technically
sound: the method is transparent and the artefacts are public; the evidence
remains narrow and is labelled as such. (3)~Comprehensive: the related work is
retained and the positioning claim is now falsifiable. (4)~and (5)~References:
applicable and sufficient; none is inappropriate.

\section{What this revision does not claim}
\label{sec:partial}

For the Associate Editor's convenience, the eight items that are met only in
part, with the single missing element in each:

\begin{center}\small
\begin{tabular}{@{}p{0.09\textwidth}p{0.42\textwidth}p{0.42\textwidth}@{}}
\toprule
\textbf{Item} & \textbf{What was done} & \textbf{What is still missing} \\
\midrule
R1.2 & Truth table over $(Y,M,O)$ with $E$ split; $A$, $D$ and $H$ defined; the
engine-versus-policy split reported & The $E{=}$unknown contract is stated across
several sections rather than collected in one place \\
R1.4 & Pre-registered protocol bound by digest; primary-vantage gate checked
before the run; amendment timing disclosed & A mechanical, implementation-
independent test of independence from the tested mechanism \\
R1.8 & Per-channel four-cell matrix for Gitea and GitLab; the two failure
mechanisms separated & The same matrix for Gogs, whose surface exposes no
equivalent decomposition \\
R1.9 / R2.4 & A real vulnerable/fixed pair through the adjudication gate; the
seeded control bounded to one fault pattern & Per-defect recall and false-positive
rates across several independent defect classes \\
R1.12 & Fixed-weight derivation; arbitrary class-share bound & A derivation for
detector-endogenous input counts \\
R1.14 & Pinned environment, commands, seeds, raw outputs, regeneration scripts,
clean-copy verification & The failure and missing-observation rules stated in the
body rather than only in the protocol file and scripts \\
R2.6 & Design dependence stated; per-block and per-cell counts inspectable &
Rates recomputed under equal per-cell weights alongside the original weights \\
\bottomrule
\end{tabular}
\end{center}

\noindent We would rather submit with these eight named than with a text that
implies they are closed. Each is visible in the manuscript as a stated boundary,
in Section~VIII or at the point of use.

\section{Files in this resubmission}

\begin{enumerate}
\item This response to reviewers.
\item The revised manuscript with all changes highlighted.
\item A clean copy of the revised manuscript, as PDF and as the LaTeX source
package (90 files), which compiles in an empty directory and reproduces every
table and figure from the raw records.
\end{enumerate}

\end{document}
